\chapter*{Note to Instructors}
\addcontentsline{toc}{chapter}{Preface and Note to Instructors}

The latest version of this book is an attempt at making this textbook
more applicable in a wider range of courses. 
Although it was originally designed for use in a follow-up course 
to a data structures course for computer science majors,
the current version can also be used for a discrete mathematics
course with few to no prerequisites, and can be more or less focused
on proof writing according the the needs of a course.

The book was 
reorganized and has had content added to make it suitable for
a discrete mathematics course with no prerequisites for 
secondary mathematics education students in the state of Michigan
while continuing to serve the computer science students. 
I will be testing it in Spring 2025 for the first time with a mixed group
of secondary mathematics education, computer science, 
and potentially general education students.

Several changes were made in the past few versions to accomplish the
goals stated above.
The content was rearranged to allow a course 
to focus on proof writing or not. 
For instance, 
Sections \ref{section:setProofs}, 
\ref{section:functionProofs}, 
\ref{subsec:alProof}, 
and
\ref{subsec:limproof}
separate proofs involving
sets, functions, and computational complexity so they can be skipped if desired.
In fact, I currently teach a version of the course that is
not proof focused and we skip the entirety of Chapter~\ref{ch:proofs} 
on proofs along with those sections, yet we are still able to cover
induction proofs and the students seem to still pick up on that technique.

I also attempted to provide more details 
about topics I previously expected were review 
(e.g. programming fundamentals, sorting and searching algorithms, matrices),
and in a few places I note 
that examples or sections can be skipped based on 
the reader's background 
(e.g. Section~\ref{subsec:BDS} on basic data structures).
In addition, a section on minimum spanning trees 
was added so that non-computer science students would see 
at least one interesting example of a graph problem and
several algorithms to solve it.

Since the textbook has gone through some major changes in the past
few years, I definitely appreciate any feedback, positive or negative
(preferably constructive, though!), so that I can continue 
improving it. So if you see things that seem out of place, 
topics that you think should be there,
notice other problems, or have suggestions, please let me know.

In the past I have used this textbook for a discrete mathematics course
for computer science students with a data structures prerequisite and 
have covered the vast majority of the textbook.
As mentioned above, it will be used for a no-prerequisite discrete mathematics course for secondary
mathematics education students in Spring 2025. The current planned coverage is as follows
\begin{itemize}
\item 
Sections \ref{section:logic}-\ref{section:pred}
\item 
Sections \ref{section:sets}-\ref{section:functions} 
(skipping \ref{section:setProofs} and 
\ref{section:functionProofs})
\item Chapter \ref{ch:algs}
\item Chapter \ref{ch:seqsummat}
\item Chapter \ref{ch:AlgAn}
\item Chapter \ref{ch:recurrences} (skipping \ref{subsec:linear})
\item Chapter \ref{ch:count} (probably skipping \ref{sec:inclExcl})
\item Chapter \ref{ch:graphs}
\end{itemize}
%I am sad to say that I have not had sufficient time to map out 
%prerequisite sections or suggest which section may be useful in
%different sorts of courses. 
%I leave it in your capable hands to determine that for yourself.\\

Here is an anecdote on how I personally use the book that you might 
find helpful.
The \verb+TL;DR+ version: I have them answers all exercises
and reading questions before class
and their unanswered questions
are where each class period begins, and often lingers.
I rarely (never?) lecture through material when I use this textbook.
Students seem to like it and they are learning.

At the beginning of each semester, I have each student 
share a Google Doc with me called 
{\em Smith RQ}, where they replace {\em Smith} with their last name.
Every time I assign reading from the textbook,
I tell them to write answers to all of the exercises in the assigned sections
of the book (my students buy a printed copy of the book), 
and to put their answers to all of the reading questions into their 
Google Doc, placing the most recent answers at the top of the document
(to make it quicker for me to find).\footnote{
This could be done the old-fashioned way on pencil and paper, 
or by having them submit it via a CMS, but using 
a Google Doc means that the deadline is flexible and
they do not have to repeatedly turn their answers in since they 
just prepend them to the document they have already shared with me.}
I quickly go through their
documents to make sure they have filled out answers to the reading 
questions for the day, 
giving them 0-4 points depending on how they did.\footnote{Why
0-4? Since 4 is even, they can get half credit. And 5 categories seems
like a good number. The way I usually assign grades, 0-2 would also work
since very rarely do I give someone 1 or 3 points. But you do you.
}
For sake
of time, I only grade them based on attempt, not content.\footnote{I
think it is very important that students know when they are being graded
on effort instead of outcome since I never want them to get a sense that
they did something correct based on the fact that I did not mark it down.
I make it very clear to my students that it is their responsibility to grade
their own work and they need to ask questions if they are uncertain
about any of their answers.
} 
If they attempted to answer all of the questions 
(including answers such as ``I'm not sure how to do this one''),
they get 4 points. For half, they get 2. If they skip it, a big fat zero!
You get the idea.
When I have time, I do this before class so I can also take a brief look at 
some of the answers to get a sense of how the class is doing on the 
material.\footnote{
When my schedule allows me a free hour before class, I have the reading
questions due an hour before class. Other semester I have them due at 
midnight the night before or mid-morning to allow me time to look at 
them before class.}


At the beginning of class, I pick a random\footnote{
By ``random'' I really mean a page that is usually
in the second half of the reading and has multiple questions on it.
}
page from the reading 
that has exercises on it and ask all of them to show me their book. 
If they attempted to answer
all of the questions on the page (both displayed pages), they get 4 points.
If they tried some of them, they get 2 or 3. If they forgot to do it, they get 0.

I count these attempts as part of their grade so that they have motivation
to do it, but since the grades {\em should}
be fairly high on these, I do not want it to count for too much.
I call this category {\em Activities} or {\em Preparation}, and
it is often worth 15-20\% of their final grade. My exams and homework
are generally sufficiently difficult to compensate for the fact that many of 
the students will have really high scores in this category.

If one had access to a grader, a nice slight modification would be to
pick a handful of questions to grade for content, and the rest for completion.
That can work for both the exercises 
(they can put select answers in their Google Doc, for instance) 
and the reading questions. On the other hand, since 
answers to the exercises and reading questions are available in the book,
I am not sure that much is gained by grading these for content.

I spend the next part of class, and often the entire class, focusing
on answering questions from the exercises and reading questions that 
stumped them. 
I rarely lecture through material in this course. 
I expect that they will learn the mundane parts
(e.g. basic definitions) through reading, and we can
focus our time during class on the more complicated concepts and applying
the ideas. 

I have used this technique for over a decade and it seems to 
work well. In addition, most of my students appreciate the way the textbook 
is written and how I use it since it forces them to do the reading 
(that they otherwise would skip doing) and they 
get useful feedback along the way. On the other hand, they do mention 
that it is a lot of work---reading, answering exercises, and answering reading
questions takes a lot longer than just skimming, which is what many of them
would likely do otherwise. But in the end, they definitely seem to fall on 
the side of liking this approach.

\bigskip

\noindent
Charles A. Cusack

\noindent
June, 2024


%\chapter*{Preface}
%\addcontentsline{toc}{chapter}{Preface and Note to Instructors}
\newpage
\noindent{\LARGE \bf Preface}

\bigskip

This book is an attempt to present some of the most important discrete mathematics concepts
to computer science students in the context of algorithms.  I wrote it for use as a textbook
for half of a course on discrete mathematics and algorithms that we offer at Hope College.
Since it was written for use in this particular context,
it is important to note that the course (and therefore this book) 
has as a prerequisite what is typically referred to as CS2.
Thus, it is assumed that the reader has had one or more programming courses
(the example code in the book is very much like C++ or Java), 
has seen basic data structures (e.g. stacks, queues, linked lists, trees), 
has studied the standard sorting algorithms (e.g. selection sort, bubble sort,
insertion sort, Quicksort, and merge sort), 
and
has seen recursion (although that material is reviewed in the book). 
In addition, it is assumed that the reader has seen the binary representation of integers.
Finally, limits and derivatives are used in a few sections in the chapter on algorithm analysis
because I (and I think many students) find it much easier to prove bounds using limits
instead of the formal definitions, assuming they have had calculus. Because the majority of
our students have had calculus, I do not review those concepts. 
I have found that even the students who have not had calculus can tackle this material 
with a little extra help, especially given the fact that the derivatives and limits they need to 
compute are pretty straightforward. 
Alternatively, that section can be skipped 
(although it makes the next section on growth rates a bit more difficult to read in detail).

Some of the material is drawn from several open-source books by David Santos.
Other material is from handouts I have written and used over the years.  
I have extensively edited the material from both sources, 
both for clarity and to emphasize the connections
between the material and algorithms where possible.
I have also added a significant amount of new material.
The format of the material is also significantly different than it was in
the original sources.

I should mention that I never met David Santos, who apparently died in 2011. 
I stumbled upon his books in the summer  of 2013 when 
I was searching for a discrete mathematics book to use 
in a new course.  
When I discovered that I could adapt his material for my own use, I decided to do so.  
Since clearly he has no knowledge of this book, 
he bears no responsibility for any of the edited content.
Any errors or omissions are therefore mine.

I appreciate any feedback you have.  
Please send any typos, formatting errors, other errors, suggestions, etc., to  
\href{mailto:cusack@hope.edu}{cusack@hope.edu}.

I would like to thank the following people for submitting feedback/errata (listed in no particular order):
Dan Zingaro, 
Mike Jipping, 
Steve Ratering, 
Victoria Gonda, 
Nathan Vance, 
Cole Watson, 
Kalli Crandell, 
John Dood, 
James Cerone,
Coty Franklin, 
Kyle Magnuson, 
Karl-Dieter Crisman,
Katie Brudos, 
Jonathan Senning, 
Matthew DeJongh, 
Julian Payne, 
Josiah Brouwer,
Cole Bruns,
and probably several others I forgot to mention (sorry!).
I would like to especially thank Lydia Won for helping 
write/compile solutions for the reading comprehension questions.

\bigskip
\noindent
Charles A. Cusack

\noindent
May, 2019

